%% kapa: exam=1 round=34 year=2023 period=1 form=A subject=civil_law subject_ko=민법 seq=9 number=9
%% source: Q-Net 감정평가사 기출문제 (cst00309) artlSeq=5211789 file=raw/1008/5211789/2023년도 제34회 감정평가사 1차시험 1교시.pdf
%% conversion: kordoc 4.24.51 -> tools/gen1.py (auto); review=proofread-sonnet; flags=none
\begin{question}{9}
법률행위의 목적에 관한 설명으로 옳은 것을 모두 고른 것은?
\begin{framedtext}
ㄱ. 甲이 乙에게 매도한 건물이 계약체결 후 甲의 방화로 전소하여 乙에게 이전할 수 없게 된 경우, 甲의 손해배상책임이 문제될 수 있다. \par
ㄴ. 甲이 乙에게 매도한 토지가 계약체결 후 재결수용으로 인하여 乙에게 이전할 수 없게 된 경우, 위험부담이 문제될 수 있다. \par
ㄷ. 甲이 乙에게 매도하기로 한 건물이 계약체결 전에 지진으로 전파(全破)된 경우, 계약체결상의 과실책임이 문제될 수 있다. \par
\end{framedtext}
\begin{choices}
  \item ㄴ
  \item ㄱ, ㄴ
  \item ㄱ, ㄷ
  \item ㄴ, ㄷ
  \item ㄱ, ㄴ, ㄷ
\end{choices}
\end{question}
\answer{5}
